% 中文审阅稿：不纳入 main.tex 编译。
% 若需单独生成 PDF，请使用 XeLaTeX 并引入 \usepackage{ctex}。
% 所有 \cite 均来自 docs/cvpr2027/parsing。

\section{引言}
\label{sec:intro_zh}

细线结构分割的本质不在于逐像素分类，而在于从图像中恢复一个保持连通的网络。对血管、神经、裂缝、道路、线缆和砌体灰缝等细线目标而言，分割质量不能仅由像素覆盖量衡量，更关键的是其网络连接关系能否被完整恢复。
然而，以像素重叠为主的评价指标与任务对连通性的实际需求并不总是一致，这导致细线分割评价中长期存在一种错位：
两个 Dice 相近的预测，可能在拓扑上一个可用、一个不可用；
而一条细分支的断裂，足以改变血流模拟、道路连通性或裂缝扩展分析的结论。
因此，对这类结构而言，连通性与拓扑并非附属指标，而是与重叠率同等重要的一阶约束 \cite{shit2021cldice,kirchhoff2024skeletonrecall}。

这类结构在真实场景中广泛存在，典型实例包括眼底与 OCTA 血管图像 \cite{staal2004drive,hoover2000stare,budai2013robust,li2024octa500}、冠状动脉 \cite{cervantes2019dca1}、路面裂缝 \cite{liu2019deepcrack}、道路 \cite{mnih2013machine}，以及线缆、轮胎花纹和砌体灰缝 \cite{zhu2026sacm,valverde2025topomortar}。
这些目标通常细长、低对比度，并常伴随大量分叉、交叉与遮挡；与此同时，稠密标注成本高昂，测试域在成像模态、分辨率与采集条件上差异显著。
上述因素共同使得少标注与跨域条件下的细线分割成为实际部署中最棘手的问题之一。



过去十年，这一领域沿着两条主线推进。
第一条是表征学习：
U-Net \cite{ronneberger2015unet} 奠定了编码器—解码器范式，
CS-Net \cite{mou2019csnet} 与 CS2-Net \cite{mou2021cs2net} 引入注意力与多尺度 设计，
DeepCrack \cite{liu2019deepcrack} 与 nnWNet \cite{zhou2025nnwnet} 分别从裂缝和 Transformer—卷积混合结构入手，不断提升局部细节与多尺度上下文。
第二条是拓扑监督：centerline Dice \cite{shit2021cldice} 与 skeleton recall \cite{kirchhoff2024skeletonrecall} 把骨架或连通性写入损失，显著缓解断裂与碎片化。
两条主线在域内都有效，但共享同一个假设：训练与测试同分布。
当模态或采集条件变化时，性能会明显下降；更重要的是，它们都把问题视为“训练一个更好的分割器”，而没有触及推理阶段的一个结构性选择——当模型给出多个候选时，如何收敛成一个输出。




SAM \cite{kirillov2023sam} 提供了另一条路径。
冻结的 ViT 编码器配合轻量适配器，使得用极少标注适配新域成为可能。
SAM-Adapter \cite{chen2023samadapter} 与 Conv-LoRA \cite{zhong2024convlora} 从 adapter 与低秩适配入手，
CAT-SAM \cite{xiao2024catsam} 做少样本条件调优，HQ-SAM \cite{ke2023hqsam} 改进高质量掩膜解码；
这些工作共同把“通用分割先验”转化为“可部署的域内分割器”。
在细线结构上，UCS \cite{zhu2023ucs} 用稀疏 adapter 与基于 FFT 的提示生成推进跨域分割；
SACM \cite{zhu2026sacm} 则用双层 adapter 与无提示解码器，仅以 18 张标注图像训练，就能在十二个数据集上取得较好的跨域泛化。
SACM 的解码器有一个值得注意的设计：它不再只输出一个掩膜，而是维护四个候选头，用 Stage-1 置信度排序、Stage-2 细化，再用一个 MLP IoU 预测器在测试时选出最终掩膜。

这引出了本文的核心问题：多候选解码的收益，究竟来自“生成更多候选”，还是来自“如何部署候选”？
现有做法把答案交给一个测试时的选择器 \cite{zhu2026sacm}。
但选择器是在训练分布上学习的，却要在未见模态上排序，它的校准会漂移，打分所依赖的“重叠”也可能偏向更粗、更连通的掩膜，牺牲细分支。
更微妙的是，选择器与候选头在同一训练流程中耦合：候选头被鼓励分化成不同角色，选择器则被鼓励挑一个在当前分布上最稳的；
在域偏移下，这种耦合可能使选择器偏向某一类候选，压掉其他角色专家。
于是，一个本应用于提升鲁棒性的多候选机制，反而把最终输出重新交给了单点决策。
我们因此提出两个问题：候选头完成角色分化之后，测试时是否还需要一个选择器？ 如果不需要，一条固定规则能否在跨域场景下替代它？



我们的出发点是：如果四个候选头已经在训练中分化成不同角色——稠密、拓扑、边界、召回——那么测试时最需要的可能不是“更聪明的选择”，而是“不选择”。
我们提出 Deterministic Candidate Decoding (DCD)：保留 SACM 的冻结 SAM ViT-L、双层 adapter 与四候选解码器；
训练时四个头都接受稠密监督与 soft clDice/skeleton recall，三个辅助头再分别接受拓扑、边界、召回三类角色损失，零初始化的线性融合仅作为辅助输出。
推理时，gating 只负责排序，DCD 始终部署固定排名 slot 3，不使用 IoU 预测器，也不使用融合分支。
slot 3 对应召回偏置的 Tversky 专家；
固定该排名既保留了训练中形成的角色分工，又移除了选择器带来的方差。
我们在同一 checkpoint 上做了对照：固定部署相对 gating+fusion 在 Dice/IoU/clDice 上分别提升 1.22/1.11/1.06，HD95 基本不变。
这个反直觉的结果说明，多候选解码的瓶颈不在“生成”，而在“部署”。



% 图 2 位置：DCD 框架图。
% 正式英文稿中对应 fig:dcd_framework。

我们在严格的 3-shot 协议下验证这一判断：训练仅使用来自六个 Seen 数据集的18 张图像，十二个评测数据集不参与训练或模型选择；
其中八个 Unseen（同一结构类别、不同数据集/模态/采集条件）与四个 Novel（训练中未出现的结构类别）。
与十个 CNN、Transformer 与 SAM 基线比较，并对四个训练种子取平均后，DCD 取得 58.76 Dice、45.17 IoU、66.02 clDice 与 63.13 HD95 的数据集宏平均，四项指标均排名第一。
我们同时报告逐数据集结果与 Dice—HD95 权衡：DCD 并非每个数据集都最好，HRF、STARE 等细血管场景仍有差距，部分数据集上 HD95 让位于重叠率。
消融进一步分离了部署规则、角色损失、拓扑监督与融合的贡献，支持“部署是瓶颈”的核心判断。

本文的贡献可以概括为三点。
\begin{itemize}
\item 我们重新审视多候选解码的收益来源，指出测试时候选选择是细线分割的部署瓶颈，并给出反直觉证据：固定排名部署在重叠与拓扑指标上优于测试时选择，且无需重新训练候选专家。
\item 我们提出 DCD：以稠密监督、拓扑感知与角色特定目标训练四个候选头，推理时只部署固定排名 slot 3 的候选；gating 只负责排序，融合分支只在训练中使用。
      它把多候选解码从输入相关的选择，转化为确定性的单输出规则。
\item 我们建立了统一的 3-shot 细线分割评测协议，覆盖 12 个数据集（8 个未见域、4 个新类）、10 个 CNN/Transformer/SAM 基线和 4 个训练种子。
      DCD 在该协议下四项指标均取得最佳宏平均。
      消融进一步表明，收益主要来自候选角色化与固定部署，而不是测试时选择或融合。
\end{itemize}
